\documentclass[10pt,a4paper]{amsart}
\usepackage{amsmath,amssymb,amsthm,graphicx}
% Retained minimal preamble for the Stacks Project chapter varieties.tex.
% Verbatim excerpt of the chapter preamble: Xy-pic setup (native preamble lines 14-18), the native theorem
% declarations with the mathematical macros (native preamble lines 154-217), and the project own no-op
% definitions of the reference/slogan/history display environments (native wrapper.tex lines 1-3) that the
% counted statement uses for its slogan. Copied mechanically; nothing added, removed or redefined.
% For commutative diagrams we use Xy-pic
\usepackage[all]{xy}

% We use 2cell for 2-commutative diagrams.
\xyoption{2cell}
% Theorem environments.
%
\theoremstyle{plain}
\newtheorem{theorem}[subsection]{Theorem}
\newtheorem{proposition}[subsection]{Proposition}
\newtheorem{lemma}[subsection]{Lemma}

\theoremstyle{definition}
\newtheorem{definition}[subsection]{Definition}
\newtheorem{example}[subsection]{Example}
\newtheorem{exercise}[subsection]{Exercise}
\newtheorem{situation}[subsection]{Situation}

\theoremstyle{remark}
\newtheorem{remark}[subsection]{Remark}
\newtheorem{remarks}[subsection]{Remarks}

\numberwithin{equation}{subsection}

% Macros
%
\def\lim{\mathop{\mathrm{lim}}\nolimits}
\def\colim{\mathop{\mathrm{colim}}\nolimits}
\def\Spec{\mathop{\mathrm{Spec}}}
\def\Hom{\mathop{\mathrm{Hom}}\nolimits}
\def\Ext{\mathop{\mathrm{Ext}}\nolimits}
\def\SheafHom{\mathop{\mathcal{H}\!\mathit{om}}\nolimits}
\def\SheafExt{\mathop{\mathcal{E}\!\mathit{xt}}\nolimits}
\def\Sch{\mathit{Sch}}
\def\Mor{\mathop{\mathrm{Mor}}\nolimits}
\def\Ob{\mathop{\mathrm{Ob}}\nolimits}
\def\Sh{\mathop{\mathit{Sh}}\nolimits}
\def\NL{\mathop{N\!L}\nolimits}
\def\CH{\mathop{\mathrm{CH}}\nolimits}
\def\proetale{{pro\text{-}\acute{e}tale}}
\def\etale{{\acute{e}tale}}
\def\QCoh{\mathit{QCoh}}
\def\Ker{\mathop{\mathrm{Ker}}}
\def\Im{\mathop{\mathrm{Im}}}
\def\Coker{\mathop{\mathrm{Coker}}}
\def\Coim{\mathop{\mathrm{Coim}}}

% Boxtimes
%
\DeclareMathSymbol{\boxtimes}{\mathbin}{AMSa}{"02}

%
% Macros for moduli stacks/spaces
%
\def\QCohstack{\mathcal{QC}\!\mathit{oh}}
\def\Cohstack{\mathcal{C}\!\mathit{oh}}
\def\Spacesstack{\mathcal{S}\!\mathit{paces}}
\def\Quotfunctor{\mathrm{Quot}}
\def\Hilbfunctor{\mathrm{Hilb}}
\def\Curvesstack{\mathcal{C}\!\mathit{urves}}
\def\Polarizedstack{\mathcal{P}\!\mathit{olarized}}
\def\Complexesstack{\mathcal{C}\!\mathit{omplexes}}
% \Pic is the operator that assigns to X its picard group, usage \Pic(X)
% \Picardstack_{X/B} denotes the Picard stack of X over B
% \Picardfunctor_{X/B} denotes the Picard functor of X over B
\def\Pic{\mathop{\mathrm{Pic}}\nolimits}
\def\Picardstack{\mathcal{P}\!\mathit{ic}}
\def\Picardfunctor{\mathrm{Pic}}
\def\Deformationcategory{\mathcal{D}\!\mathit{ef}}
\newenvironment{reference}{}{}
\newenvironment{slogan}{}{}
\newenvironment{history}{}{}
\usepackage{cleveref}
\crefname{originalreference}{original reference}{original references}
\Crefname{originalreference}{Original reference}{Original references}
\makeatletter
\newlabel{algebra-lemma-support-point}{{00L5}{source}{}{sourcepointer-1}{}}
\newlabel{algebra-lemma-support-point@cref}{{[originalreference][0][]00L5}{[1][0][]source}{}{}{}}
\makeatother
\crefname{originallemma}{original lemma}{original lemmas}
\Crefname{originallemma}{Original lemma}{Original lemmas}
\makeatletter
\newlabel{lemma-euler-characteristic-additive}{{08AA}{source}{}{sourcepointer-2}{}}
\newlabel{lemma-euler-characteristic-additive@cref}{{[originallemma][0][]08AA}{[1][0][]source}{}{}{}}
\makeatother
\crefname{originalreference}{original reference}{original references}
\Crefname{originalreference}{Original reference}{Original references}
\makeatletter
\newlabel{algebra-lemma-length-additive}{{00IV}{source}{}{sourcepointer-3}{}}
\newlabel{algebra-lemma-length-additive@cref}{{[originalreference][0][]00IV}{[1][0][]source}{}{}{}}
\makeatother
\crefname{originalreference}{original reference}{original references}
\Crefname{originalreference}{Original reference}{Original references}
\makeatletter
\newlabel{coherent-lemma-coherent-filter}{{01YF}{source}{}{sourcepointer-4}{}}
\newlabel{coherent-lemma-coherent-filter@cref}{{[originalreference][0][]01YF}{[1][0][]source}{}{}{}}
\makeatother
\crefname{originallemma}{original lemma}{original lemmas}
\Crefname{originallemma}{Original lemma}{Original lemmas}
\makeatletter
\newlabel{lemma-numerical-polynomial-from-euler}{{0BEM}{source}{}{sourcepointer-5}{}}
\newlabel{lemma-numerical-polynomial-from-euler@cref}{{[originallemma][0][]0BEM}{[1][0][]source}{}{}{}}
\makeatother
\crefname{originalreference}{original reference}{original references}
\Crefname{originalreference}{Original reference}{Original references}
\makeatletter
\newlabel{cohomology-section-projection-formula}{{01E6}{source}{}{sourcepointer-6}{}}
\newlabel{cohomology-section-projection-formula@cref}{{[originalreference][0][]01E6}{[1][0][]source}{}{}{}}
\makeatother
\crefname{originalreference}{original reference}{original references}
\Crefname{originalreference}{Original reference}{Original references}
\makeatletter
\newlabel{coherent-lemma-relative-affine-cohomology}{{089W}{source}{}{sourcepointer-7}{}}
\newlabel{coherent-lemma-relative-affine-cohomology@cref}{{[originalreference][0][]089W}{[1][0][]source}{}{}{}}
\makeatother
\title{The leading term of the Euler characteristic of a coherent module twisted by invertible sheaves is given by the lengths at the generic points of the top-dimensional components}
\author{The Stacks Project Authors}
\date{stacks-project:varieties.tex:local-snapshot-20261011; GNU FDL 1.2 or later}
\begin{document}
\maketitle
\noindent\textit{Editorial note.} All delivered mathematical text is the author's own, copied verbatim from the pinned Stacks Project chapter file varieties.tex; no statement or proof step was rewritten, reordered or removed.\par\medskip
\noindent\textit{Editorial note.} The counted span cites seven cross-chapter labels (including a section), each recorded in assembly.external\_references with its canonical tag as a navigation pointer.\par\medskip
\noindent\textit{Editorial note.} The proof cites no bibliography item inside the counted span.\par\medskip
\noindent\textit{Allowed standard prerequisites (editorial).}\begin{itemize}
\item Filtration of a coherent module with subquotients pushed forward from integral subschemes. Stacks Project, Cohomology of Schemes, Lemma coherent-lemma-coherent-filter (tag 01YF). Supplies the filtration that reduces the lemma to pushforwards of ideal sheaves.
\item Additivity of the Euler characteristic and of lengths. Stacks Project, Varieties, Lemma lemma-euler-characteristic-additive (tag 08AA) and Algebra, Lemma algebra-lemma-length-additive (tag 00IV). Gives additivity of the functional E along short exact sequences.
\item Numerical polynomial of the Euler characteristic and behaviour under pushforward. Stacks Project, Varieties, Lemma lemma-numerical-polynomial-from-euler (tag 0BEM); Cohomology of Schemes, Lemma coherent-lemma-relative-affine-cohomology (tag 089W). Controls the degree of the Euler characteristic polynomial and identifies the cohomology of a closed-immersion pushforward.
\end{itemize}\medskip
\ifdefined\hypertarget\hypertarget{sourcepointer-1}{}\fi
\noindent\textit{Source reference.} Original reference 00L5: sources/stacks-chow/tags.txt, source lines 775--775. Cross-chapter standard prerequisite cited by the author inside the counted span; it is used as a pointer and its own proof is not part of this item.\par\medskip
\ifdefined\hypertarget\hypertarget{sourcepointer-2}{}\fi
\noindent\textit{Source reference.} Original lemma 08AA: sources/stacks-chow/tags.txt, source lines 10525--10525. Cross-chapter standard prerequisite cited by the author inside the counted span; it is used as a pointer and its own proof is not part of this item.\par\medskip
\ifdefined\hypertarget\hypertarget{sourcepointer-3}{}\fi
\noindent\textit{Source reference.} Original reference 00IV: sources/stacks-chow/tags.txt, source lines 693--693. Cross-chapter standard prerequisite cited by the author inside the counted span; it is used as a pointer and its own proof is not part of this item.\par\medskip
\ifdefined\hypertarget\hypertarget{sourcepointer-4}{}\fi
\noindent\textit{Source reference.} Original reference 01YF: sources/stacks-chow/tags.txt, source lines 2554--2554. Cross-chapter standard prerequisite cited by the author inside the counted span; it is used as a pointer and its own proof is not part of this item.\par\medskip
\ifdefined\hypertarget\hypertarget{sourcepointer-5}{}\fi
\noindent\textit{Source reference.} Original lemma 0BEM: sources/stacks-chow/tags.txt, source lines 14355--14355. Cross-chapter standard prerequisite cited by the author inside the counted span; it is used as a pointer and its own proof is not part of this item.\par\medskip
\ifdefined\hypertarget\hypertarget{sourcepointer-6}{}\fi
\noindent\textit{Source reference.} Original reference 01E6: sources/stacks-chow/tags.txt, source lines 1823--1823. Cross-chapter standard prerequisite cited by the author inside the counted span; it is used as a pointer and its own proof is not part of this item.\par\medskip
\ifdefined\hypertarget\hypertarget{sourcepointer-7}{}\fi
\noindent\textit{Source reference.} Original reference 089W: sources/stacks-chow/tags.txt, source lines 10511--10511. Cross-chapter standard prerequisite cited by the author inside the counted span; it is used as a pointer and its own proof is not part of this item.\par\medskip
\begin{lemma}
\label{lemma-numerical-polynomial-leading-term}
Let $k$ be a field. Let $X$ be a proper scheme over $k$. Let
$\mathcal{F}$ be a coherent $\mathcal{O}_X$-module. Let
$\mathcal{L}_1, \ldots, \mathcal{L}_r$ be invertible $\mathcal{O}_X$-modules.
Let $d = \dim(\text{Supp}(\mathcal{F}))$.
Let $Z_i \subset X$ be the irreducible components
of $\text{Supp}(\mathcal{F})$ of dimension $d$. Let $\xi_i \in Z_i$
be the generic point and set
$m_i = \text{length}_{\mathcal{O}_{X, \xi_i}}(\mathcal{F}_{\xi_i})$.
Then
$$
\chi(X, \mathcal{F} \otimes \mathcal{L}_1^{\otimes n_1} \otimes \ldots \otimes
\mathcal{L}_r^{\otimes n_r}) -
\sum\nolimits_i
m_i\ \chi(Z_i, \mathcal{L}_1^{\otimes n_1} \otimes \ldots \otimes
\mathcal{L}_r^{\otimes n_r}|_{Z_i})
$$
is a numerical polynomial in $n_1, \ldots, n_r$ of total degree $< d$.
\end{lemma}

\begin{proof}
Consider pairs $(\xi , Z)$ where $Z \subset X$ is an integral
closed subscheme of dimension $d$ and $\xi$ is its generic point.
Then the finite $\mathcal{O}_{X, \xi}$-module $\mathcal{F}_\xi$
has support contained in $\{\xi\}$ hence the length
$m_Z = \text{length}_{\mathcal{O}_{X, \xi}}(\mathcal{F}_\xi)$
is finite (Algebra, Lemma \ref{algebra-lemma-support-point})
and zero unless $Z = Z_i$ for some $i$. Thus the expression
of the lemma can be written as
$$
E(\mathcal{F}) =
\chi(X, \mathcal{F} \otimes \mathcal{L}_1^{\otimes n_1} \otimes \ldots \otimes
\mathcal{L}_r^{\otimes n_r}) -
\sum\nolimits
m_Z\ \chi(Z, \mathcal{L}_1^{\otimes n_1} \otimes \ldots \otimes
\mathcal{L}_r^{\otimes n_r}|_Z)
$$
where the sum is over integral closed subschemes $Z \subset X$
of dimension $d$. The assignment $\mathcal{F} \mapsto E(\mathcal{F})$
is additive in short exact sequences
$0 \to \mathcal{F} \to \mathcal{F}' \to \mathcal{F}'' \to 0$
of coherent $\mathcal{O}_X$-modules whose support has dimension
$\leq d$. This follows from additivity of Euler characteristics
(Lemma \ref{lemma-euler-characteristic-additive})
and additivity of lengths
(Algebra, Lemma \ref{algebra-lemma-length-additive}).
Let us apply Cohomology of Schemes, Lemma \ref{coherent-lemma-coherent-filter}
to find a filtration
$$
0 = \mathcal{F}_0 \subset \mathcal{F}_1 \subset
\ldots \subset \mathcal{F}_m = \mathcal{F}
$$
by coherent subsheaves such that for each $j = 1, \ldots, m$
there exists an integral closed subscheme $V_j \subset X$
and a nonzero sheaf of ideals $\mathcal{I}_j \subset \mathcal{O}_{V_j}$
such that
$$
\mathcal{F}_j/\mathcal{F}_{j - 1}
\cong (V_j \to X)_* \mathcal{I}_j
$$
It follows that $V_j \subset \text{Supp}(\mathcal{F})$ and
hence $\dim(V_j) \leq d$.
By the additivity we remarked upon above it suffices to
prove the result for each of the subquotients
$\mathcal{F}_j/\mathcal{F}_{j - 1}$. Thus it suffices to prove
the result when $\mathcal{F} = (V \to X)_*\mathcal{I}$ where
$V \subset X$ is an integral closed subscheme of dimension $\leq d$
and $\mathcal{I} \subset \mathcal{O}_V$ is a nonzero coherent
sheaf of ideals. If $\dim(V) < d$ and more generally for $\mathcal{F}$
whose support has dimension $< d$, then the first term
in $E(\mathcal{F})$ has total degree $< d$ by
Lemma \ref{lemma-numerical-polynomial-from-euler}
and the second term is zero. If $\dim(V) = d$, then we can use the
short exact sequence
$$
0 \to (V \to X)_*\mathcal{I} \to (V \to X)_*\mathcal{O}_V
\to (V \to X)_*(\mathcal{O}_V/\mathcal{I}) \to 0
$$
The result holds for the middle sheaf because
the only $Z$ occurring in the sum is $Z = V$
with $m_Z = 1$ and because
$$
H^i(X, ((V \to X)_*\mathcal{O}_V) \otimes 
 \mathcal{L}_1^{\otimes n_1} \otimes \ldots \otimes
\mathcal{L}_r^{\otimes n_r}) =
H^i(V,  \mathcal{L}_1^{\otimes n_1} \otimes \ldots \otimes
\mathcal{L}_r^{\otimes n_r}|_V)
$$
by the projection formula
(Cohomology, Section \ref{cohomology-section-projection-formula}) and
Cohomology of Schemes, Lemma
\ref{coherent-lemma-relative-affine-cohomology};
so in this case we actually have $E(\mathcal{F}) = 0$.
The result holds for the sheaf on the right because its support
has dimension $< d$. Thus the result holds for the sheaf on the
left and the lemma is proved.
\end{proof}

\end{document}
